\section{Introduction}
\label{sec:intro}

Multi-agent systems built on large language models (LLMs) are now deployed for
consequential reasoning tasks: medical diagnosis support, financial risk assessment,
and legal document analysis.
A central design question for these systems is \emph{which models to include}.
Most production deployments use homogeneous ensembles---multiple instances of the same model---because
homogeneous agents are easy to orchestrate and their behavior is predictable.
But homogeneous agents share training data, fine-tuning objectives, and safety alignment,
which means their errors are likely correlated.
When agents fail for the same reasons, ensemble aggregation provides no benefit.

Classical ensemble learning theory offers a clear prescription: diversity reduces correlated
errors and improves robustness~\citep{wu2023heterogeneous}.
For LLMs, the question is whether \emph{capability-level diversity}---mixing small
and large models within the same family---provides benefits analogous to those
seen in classical heterogeneous ensembles of neural networks.
Smaller models rely more on surface pattern matching while larger models apply
deeper multi-step reasoning, suggesting different failure modes.
Yet both are trained by the same provider, share safety alignment,
and are exposed to similar pre-training distributions.

{\bf What is the effect of within-family capability diversity on collective reasoning?}
Prior work on multi-agent LLM debate~\citep{du2023debate,liang2023mad}
uses homogeneous agents and does not study error correlation.
\citet{chen2023reconcile} demonstrate that heterogeneous model ensembles outperform
homogeneous ones, but use models from different providers (GPT, PaLM, Claude),
making it unclear whether capability-level diversity within a single family suffices.
\citet{rosales2025diverse} challenge the naive diversity hypothesis, showing that
question interpretation diversity often outperforms model diversity---and argue
that what matters is \emph{complementary failure modes}, not independence alone.
No prior work systematically compares capability-level diversity with
adversarially-designed tasks that specifically target shared failure modes,
while measuring pairwise error correlation.

We fill this gap with a controlled empirical study.
We construct \advbench, a 30-question adversarial benchmark covering
five cognitive trap categories, and evaluate five conditions spanning a
diversity continuum: single agents, homogeneous three-agent voting,
heterogeneous two-small plus one-large voting, and a structured debate protocol.
We use \haiku (claude-haiku-4-5) as the small model
and \sonnet (claude-sonnet-4-6) as the large model,
both from Anthropic's Claude family, enabling a clean
within-family capability diversity comparison.

Our main findings are:
(1) both models achieve high accuracy (93.3\%) individually,
leaving little headroom for ensemble improvement;
(2) the cross-capability pair has lower error correlation ($r{=}0.464$)
than within-capability pairs ($r{=}0.695$), confirming theoretical predictions;
(3) ensemble methods improve accuracy by 3.4 percentage points to 96.7\%,
but this improvement is not statistically significant ($p{=}1.000$);
(4) the debate protocol \emph{degrades} performance due to sycophantic revision;
and (5) the dominant shared failure is a safety-aligned refusal shared
identically across all conditions, which no ensemble can overcome.

Our contributions are:

\begin{itemize}[leftmargin=*,itemsep=0pt,topsep=0pt]
    \item We construct \advbench, a 30-question adversarial reasoning benchmark
          across five cognitive trap categories, with verified ground-truth answers
          and six paraphrase variants for robustness testing.
    \item We conduct a controlled empirical comparison of five ensemble conditions
          spanning a diversity continuum, using real API calls to two Claude models,
          measuring accuracy, error correlation, and perturbation robustness.
    \item We quantify that within-family capability diversity produces lower
          error correlation ($r{=}0.464$ vs.\ $r{=}0.695$), supporting the
          diversity hypothesis at the correlation level while failing to produce
          measurable accuracy gains on a near-saturated benchmark.
    \item We identify safety alignment as a \emph{cross-cutting shared bias}
          that is orthogonal to capability-level diversity and cannot be
          addressed by within-family ensemble design.
\end{itemize}

We describe related work in \secref{sec:related}, our methodology in \secref{sec:method},
results in \secref{sec:results}, discussion in \secref{sec:discussion},
and conclude in \secref{sec:conclusion}.
